\section{Introduction}
\label{sec:intro}

A language model used over a long period has to learn new tasks without losing old ones. With
low-rank adapters \citep{hu2022lora} each task adds a small update that can be kept unchanged for
ever, and what remains is which adapters should act on a given input; the model is not told the
task. If all adapters are left on, each new one changes what the model computes on earlier inputs,
and much of the old tasks is lost (Section~\ref{sec:results}: with five tasks, plain
adapters keep \FivePlainAP{} of \FivePlainLearned{} points). Published methods either restrict what a
new adapter may learn \citep{wang2023olora,liang2024inflora,spacelora2026} or add modules that weight
the adapters per input \citep{liang2025gainlora,zhao2024sapt,latentlora2026}. Both reduce forgetting.
To our knowledge none states that, for a given input, the answer \emph{will not} change when the
next task is learned.

The principle of this paper is simple. We give every task one frozen set of adapters, and a switch
sends each input to one task; we call this system FK-LoRA, after the rule \emph{finders keepers}. An answer can then change only if its route changes
(Lemma~\ref{lem:route}), so retention is a property of the rule that sets the route, and it can be
proved or measured for that rule.

The instruction gives the task away. Benchmarks hide this, because their inputs begin with the
instruction of the task. Without it, over three seeds, the two published methods that we ran change
\AudSeedsGainloraChanged{} (GainLoRA) and \AudSeedsInfloraChanged{} (InfLoRA) of \AudSeedsCompared{}
answers, against \FiveGainChanged{} and \FiveInfChanged{} with it (Section~\ref{sec:noinst}). No
benchmark input occurs in two tasks, and \recog{} below shows that this loss was not forced. What is forced is that every method
without instructions loses at least the share of inputs that two tasks have in common with
different right answers (Proposition~\ref{prop:noinstructionprice}, Appendix~\ref{sec:audit}).

Our main rule, \recog{}, reads no instruction; under its variant \handicap{} an older task keeps what
it found. \recog{} is a recogniser on frozen states of the model itself,
the same for every order of arrival, and it blocks no task. With it come an audit, which estimates
from held-out inputs the share of answers that changed (Proposition~\ref{prop:audit}), and a
handicap, a bound known in advance on what a later task can take (Theorem~\ref{thm:handicap}). On 15
benchmark tasks, which were development data, it gives \AudDevLLamOneEnd{}\% of the test inputs to
their own task, \AudScaleLongThreeEnd{}\% on 15 tasks of a second benchmark, \AudFreshFourthLEnd{}\%
on 30 unseen tasks and \AudScaleAllHundredTwentyEnd{}\% with 120 tasks (routes only). Trained
without instructions on one sequence of five of these tasks (three seeds), plain adapters behind it
end at \AudSeedsRidgeAPOne{} Rouge-L against \AudSeedsGainloraAPOne{} and \AudSeedsInfloraAPOne{},
and change \AudSeedsRidgeChanged{} of \AudSeedsCompared{} answers. This beats the published methods in
this setting, by a line fixed before the run. With a tenth of their epochs most of that gap is
learning; with their full epochs (one seed) the published methods learn the tasks and then forget
\AudFullGainloraFTInt{} and \AudFullInfloraFTInt{} points, ending at \AudFullGainloraAPOne{} and
\AudFullInfloraAPOne{}, against \AudFullOursRidgeAPOne{} for ours with the same epochs. The five tasks are among the 15 on which \recog{} was chosen.

When inputs do carry an instruction, an exact statement is possible. The idea is the rule children
use for things found on the playground: \emph{finders keepers}. Each task claims a protected zone
around the frozen word-vectors of its instruction when it is learned; nothing is trained for it, and
it is never changed. An input inside a zone belongs to the oldest task that claims it, for ever, and
its answer cannot change; it comes with a certificate. We call this rule \cert{}. The price is paid in the open: a new task
whose inputs fall inside an older zone cannot be learned on those inputs, a certified answer can
never improve, and only inputs inside a zone are covered (on the benchmark the original and light
wordings; on machine paraphrases of real tasks fewer than half). Among \ManyTasks{} real tasks about
one in six collided with an older task (Section~\ref{sec:many}). The routes of \recog{} do not
depend on the wording: with instructions and heavy rewording it ends at \AudUThreeHeavyRidgeAPOne{}
where the zones end at \AudUThreeHeavyCertAPOne{} (three seeds; Section~\ref{sec:fromexample}).

\paragraph{Contributions.}
\begin{itemize}
  \item \textbf{The principle and the price.} With one frozen set of adapters per task, an answer
  changes only if its route changes (Lemma~\ref{lem:route}); without instructions every method
  loses at least the conflict of the inputs of two tasks (Proposition~\ref{prop:noinstructionprice}).
  \item \textbf{\recog{}, an audit and a handicap.} A rule without instructions that blocks no task,
  with an estimate of the share of changed answers from held-out inputs and a bound in advance,
  which is loose (Proposition~\ref{prop:audit}, Theorem~\ref{thm:handicap}).
  \item \textbf{The published methods without instructions, and the comparison.} Both hold their
  answers less well, also with their full epochs; plain adapters behind \recog{} beat them in this setting, by a line fixed
  before the run, with the caveats above.
  \item \textbf{The certified switch, its proof and its cost at scale.} The answer to every certified
  input is the same after any number of later tasks (Theorem~\ref{thm:writeonce}); \FiveCertChanged{}
  of \FiveCertCompared{} answers changed on five tasks. About one real task in six collides with an older one, and on
  machine paraphrases the zones certify \WidenRefTfivePLow{}\% to \WidenRefTfivePHigh{}\% of the inputs
  for their own task.
  \item \textbf{An evaluation with lines fixed before the runs,} in a reduced setting (T5-base,
  mostly a tenth of the published epochs, free GPUs); changed answers are counted, and learning and
  forgetting are reported separately.
\end{itemize}

\paragraph{What we do not claim.} Not that forgetting is solved: an input outside
all zones has no guarantee, and the audit and the certificate bound change, not quality. Not the published setting: all GPU experiments use
T5-base instead of the T5-large of the GainLoRA paper (whose setting we use), and all but a few a shorter budget. Not that published methods only look good because of the
instruction. Not
that the recogniser is new: it is a known kind of classifier \citep{mcdonnell2023ranpac,zhuang2022acil},
and recognising a task from frozen word-vectors is not new \citep{latentlora2026}.
